\chapter{Unsupervised transcriptome analysis}

\paragraph{Overview: variability in biological systems}

In all the sections described before, we applied the same biological model to the patients of a given cohort. However, we observe strong variability between the samples, even within the same individual. The variability occurs at three levels: at the environmental level (disease state, tissue location, . . . ), the genotypical level (presence of individual mutations inducing varying transcriptomic activity) and even at the cell population level. 

In the next sections, we will first detail canonical clustering methods, namely \Glspl{gmm} that are still relevant in the exploratory stage to unravel homogeneous groups in a seemingly highly variable disease condition, both at the phenotypical and transcriptomic level. Then, we review some of the most common deconvolution methods designed to measure and for some of them, address the bias induced by the strong variability of the cell composition.



\section{Mixture models} 
\paragraph{Overview: introduction to mixture models}

\subsection{Gaussian mixture models}

\subsection{Introduction to the EM algorithm}
\subsubsection{Motivation}
\subsubsection{Derivation of the EM algorithm in \glsentryshort {gmm}}
\subsubsection{Variants of the EM algorithm}


\subsubsection{Application of the EM algorithm to censored data (or Poisson truncated distributions)}


\subsection{Statistical results with GMMs}

\subsubsection{Model selection}

\subsubsection{Confidence intervals}
\todo[fancyline]{alexadran thesis about derivation via the fisher informaiton matrix}




\section{State-of-the-Art Approaches for Cellular deconvolution}
\paragraph{Overview: cellular deconvolution, a promising but challenging method to unravel complex biological structures}

\subsection{Physical approaches}


\subsection{Computational approaches}

\subsubsection{Reference-free approaches}

\subsubsection{Marker-based approaches}

\subsubsection{Partial-based approaches}



